\scanpage{211}
\mspage{108}

de respecter le produit scalaire \unsure{donc|dans} la relation d'incidence $R$, mais l'inconvénient que les applications ne commutent plus à l'op de $\Lambda^{\unsure{\alpha|\times|*}}$ sur \deleted{$\unsure{M}$} $S$ et $S^\vee$, mais qu'il \deleted{faut} écrire $\varphi_0(\lambda s)=\lambda^{\unsure{\vee-1|-1}}\varphi_0(s)$ (au lieu de $\lambda\varphi_0(s)$). L'avantage d'avoir une loi de groupe bien en évidence sur les homothéties, dans cette notation, peut paraître prépondérant. Mais il faut bien noter qu'en général, \uuline{on n'a pas} d'opération du groupe $Z$ sur l'ens $\Pi$ des ``jauges'' du prégraphe $\Gamma$.

\deleted{On a bien défini une application}

\[
\begin{aligned}
&\deleted{Z_{\unsure{s|S}}\times\Pi\longrightarrow\Pi}\\
&\deleted{(\lambda,\varphi_0)\longmapsto\varphi_0\lambda^{-1}=\lambda^\vee\varphi_0,}
\end{aligned}
\]

Pour que $(\lambda,\lambda^\vee)\in\unsure{\mathfrak S|\mathscr S}_S\times\mathfrak S_{S^\vee}$ définisse \deleted{une} -- à partir de \deleted{\unsure{à}} une jauge $\varphi_0\in\Pi$ -- une deuxième jauge $\varphi$ \inserted{$=\lambda^\vee\varphi_0=\varphi_0\lambda^{-1}$}, il faut en effet, en plus de la condition que $(\lambda,\lambda^\vee)$ soit un autom du prégraphe $\Gamma$, la condition (\uline{qui dépend du choix de $\varphi_0$})

\[
\lambda^\vee\varphi_0=\varphi_0\lambda^{-1}\quad\text{i.e.}\quad\varphi_0=\lambda^\vee\varphi_0\lambda. \tag*{(65)}
\]

Si on décide d'oublier $\lambda^\vee$ (p. ex.) pour ne retenir que $\lambda$ (et définir $\varphi$ par $\varphi_0\lambda^{-1}$, \inserted{et $\dot\lambda^\vee$ par $\lambda^\vee=\varphi_0^{-1}(\lambda^{-1})$ \deleted{\unsure{p.}}}), on est confronté à ce \deleted{\unsure{spectacle}}\unsure{problème|première}, que la condition sur $\lambda$ qui exprime le \deleted{symétrie} fait que

\[
(\lambda,\lambda^\vee)=(\lambda,\varphi_0^{-1}(\lambda^{-1}))
\]

\deleted{satisfasse à \unsure{une} la condition de symétrie, (} invarie $R$ \unsure{1)|1\textsuperscript{o}} dépend de $\varphi_0$, et 2\textsuperscript{o}) ne définit pas à priori un \uline{sous-groupe} de $\mathfrak S_S$ ! En fait, si on exprime

% archive sheet number: 210
